\chapter{Hydrogen Breath Test}

\section*{What Is This Test?}
A hydrogen breath test measures hydrogen, and sometimes methane, in exhaled breath after the patient drinks a specified carbohydrate solution. Gut bacteria produce these gases when they ferment carbohydrate that is not absorbed normally or reaches the small intestine in excess. The test can help evaluate lactose or fructose malabsorption and, with an appropriate protocol, suspected small-intestinal bacterial overgrowth (SIBO).

\section*{Alternative Names}
H$_2$ Breath Test; Hydrogen-Methane Breath Test; Lactose Breath Test; Fructose Breath Test; SIBO Breath Test; Glucose Breath Test; Lactulose Breath Test

\section*{Why It Is Ordered}
Investigation of bloating, abdominal pain, diarrhea, gas, or symptoms reproducibly triggered by particular carbohydrates; evaluation of suspected lactose or fructose malabsorption; and selected assessment for SIBO. The substrate and interpretation must match the clinical question.

\section*{How This Test Is Performed}
After a baseline breath sample, the patient drinks a measured solution of lactose, fructose, glucose, or lactulose. Breath samples are collected at regular intervals for several hours while symptoms are recorded. A rise in hydrogen or methane is interpreted using the laboratory's validated protocol and the patient's symptoms; test cutoffs and accuracy vary by substrate and method.

\section*{Preparation, Risks, and Limitations}
Preparation commonly includes a restricted diet the day before, fasting overnight, avoiding smoking and vigorous exercise during testing, and temporarily changing selected medicines only as directed. Antibiotics, laxatives, probiotics, bowel preparation, and recent gastrointestinal infection can affect results. The drink may cause temporary bloating, cramps, or diarrhea. A positive result does not by itself establish the cause of symptoms, and false positives and false negatives occur.

\section*{References}
\begin{enumerate}
  \item MedlinePlus. Lactose Tolerance Tests.
  \item North American Consensus on Breath Testing in Gastrointestinal Disorders.
\end{enumerate}

\clearpage
\section*{Regulatory Status and Authorization Date}
Regulatory status applies to the specific assay, instrument, manufacturer, and intended use rather than to the test name alone. No single approval date can be assigned to this test category. For a commercial assay, verify the current product in the relevant regulator's database: the U.S. Food and Drug Administration (FDA) clearance, approval, or authorization; India's Central Drugs Standard Control Organisation (CDSCO) licence or permission; the European Union In Vitro Diagnostic Regulation (EU IVDR) CE certificate issued through the applicable conformity-assessment route; or the United Kingdom Medicines and Healthcare products Regulatory Agency (MHRA) registration and UKCA/CE status. Laboratory-developed tests may not have an individual FDA, CDSCO, EU IVDR, or MHRA approval date.
\clearpage
